\documentclass[11pt]{article}
\usepackage{hyperref}
\title{A new merge claim can block an unchanged handoff}
\author{}
\date{}
\begin{document}
\maketitle

A warehouse and a railway replay use the same four movements at the same coordinates. The warehouse moves a three-carrier chain first. The railway can move a lone competing train first, depending on its native merge priority. All three movements on the original chain can still be authorized and physically possible in isolation. A proof that tracks only those movements therefore cannot certify the handoff when another actor adds a claim to their shared exit.

\paragraph{Plant and service.}
The old holder O must leave the port while retaining an earlier duty. New carrier N can use the port as O leaves. O needs X to move first, and X needs the empty merge cell. A fourth actor Q also has a valid approach to the merge by a different edge. Q owes a move into that cell. These two requests have different service effects: N needs an actual delivery or arrival at the port; Q needs actual use of the merge. Only one actor can reach an exclusive merge cell at a cut. The winner can leave it on the next cut if its following move remains legal. A cut compares placement after native arbitration, not whether an intent was sent. Once a competing grant is issued, a receiver may not silently revoke it to rescue the old proof.

The warehouse uses selected native RWARE methods pinned by \href{https://github.com/eternal-november/eos-research-programmes/pull/429}{the two-source receiver study}. The graph shim extends the earlier native harness exactly on acyclic components with a unique longest path: it enumerates the paths and refuses an ambiguous maximum. The rail fixture executes the selected Flatland candidate, motion-check, distribute and arrival methods from the same pin as that study. Fixed candidate geometry and boundary flags, speed sink and irrelevant reward sink are supplied as in that study; this is not the full RailEnv. Some grant omissions are modeled by holding the corresponding train in its old block, not by asserting that a network event was captured. The effect-side receiver file from the earlier study is unchanged; its source-specific data still include live old return or trip.

\paragraph{The matched rule.}
For this merge shape only, the old branch is a downstream suffix of the three requested moves: X must be allowed, then O can follow, then N can follow O. Without Q's claim it succeeds to that extent. With both claims admitted, RWARE uses the strict longest chain; Flatland's motion checker resolves its contested target by handle priority, and losing leading moves stop their followers. If Q is entitled to move and outranks X, introducing Q alone moves the earliest possible N service from tick 1 to tick 2 in the paired rail controls. If Q is lower priority, the same added intent does not change the outcome. The source change is therefore in how another actor's admitted claim can affect the old branch, not in the physical length of that branch.

For the all-grants control, RWARE and X-first Flatland deliver N at 1 and move Q at 2. Q-first Flatland moves Q at 1 and delivers N at 2. Both old obligations remain owed; winning the port transfer does not end them. Withholding Q's claim \emph{before} issue restores the N-first schedule even when Q would have won the native arbitration. Withholding X before issue lets Q go first in the warehouse. These are lawful alternatives only if the applicable deadline and priority authority allow delaying the other job. If both jobs are due at tick 1 neither order can satisfy both. If it is not yet known which job has the tick-1 obligation, and no lawful way to obtain the choice before issue exists, no fixed order meets the two possible laws at that cut. A successful second tick also needs the winner actually able to leave. If the rail configuration instead makes Q stop in the merge at its terminal target and retain that block, N cannot use the same repair; replay confirms no N service through tick 2.

\paragraph{What a proof must depend on.}
In the distinct-target case of \href{https://github.com/eternal-november/eos-research-programmes/pull/428}{the earlier movement study}, a fixed set of jointly required grants can suffice. At the merge a sufficient interface also bounds competing claims and the priority rule at the governed cut; in this one-competitor model, the presence of a higher-priority claim is an exact additional dependency. Adding that claim is a counterexample to monotone reuse even with the successful branch's route, cargo and every one of its grants unchanged. The claim that the winner releases the block by the next cut is separate and requires real next-step premises. A receiver reading only the completed transfer does not need to know either scheduler's priority rule, but promising the transfer before it happens does.

The direct comparison computes the strict path maximum or handle winner on the same input and ties. It agreed with actual execution for all 398 first-cut controls (14 non-tied warehouse configurations and 384 railway grant and priority combinations); both all-grant two-cut schedules and the retained-terminal counterexample execute. This is a checked limitation on proof reuse and a native cross-source distinction. There is no advantage over exact matched arbitration, claim about priority ties or unrestricted native graph planning, independently chosen workload, or R2/I1 acceptance.
\end{document}
